% 中文摘要頁
\begin{ZhAbstract}
    \begin{ZhAbstractItems}
        % 關鍵詞，請自己填，請自己填，多個關鍵字以逗號(、)隔開
        \noindent \text 關鍵詞：大型語言模型代理人、API 任務自動化、原子化經驗、動態局部修復

    \end{ZhAbstractItems}

    \begin{ZhAbstractDescription}
        大型語言模型代理人能透過應用程式介面（Application Programming Interface, API）操作外部系統，但在多步驟任務中，若規劃不完整、資料依賴處理錯誤，或個別步驟執行失敗，後續步驟可能因缺少正確的中間結果而無法完成，進而影響整體任務成果。本研究採用的基礎系統為 API 知識圖譜代理人（API Knowledge Graph Agent, APIKGAgent），其主要透過重新規劃處理執行失敗，尚未整合由成功任務軌跡建立的經驗重用機制與局部計畫修復。因此，本研究探討如何運用既有成功經驗，並針對局部錯誤調整計畫，以提升代理人的任務完成能力。
        本研究以 APIKGAgent 為基礎，提出案例式修復 API 知識圖譜代理人（Case-Based Repair API Knowledge Graph Agent, CBR-APIKGAgent），整合原子化經驗（Atomic Experience）重用與動態局部修復（Dynamic Local Repair）。系統從訓練資料的成功任務軌跡中萃取可跨任務重用的策略單元，並依當前任務或失敗情境檢索相關經驗，提供初始規劃、局部修復與重新規劃參考。當步驟執行失敗時，系統根據錯誤訊息、相依資料與執行結果，透過插入、替換或刪除步驟調整局部計畫；若無法形成有效修正，再回退至重新規劃。
        本研究於 AppWorld 的 168 個多應用程式任務上進行評估。完整系統的任務目標完成率由基礎系統的 37.5\% 提升至 51.8\%，情境目標完成率則由 14.3\% 提升至 35.7\%。消融實驗與案例分析顯示，原子化經驗主要協助代理人掌握資料依賴、目標集合與操作條件，動態局部修復則能處理授權缺失、前置資料不足及局部操作不完整等執行錯誤；兩項機制分別作用於規劃與執行階段，結合後取得最佳整體表現。研究亦發現，系統表現仍會受到檢索經驗與任務限制的適配程度、初始規劃品質，以及局部修復與重新規劃判斷等因素影響。儘管存在上述限制，整體結果仍顯示，細粒度經驗重用與局部計畫修復有助於提升大型語言模型代理人在多步驟 API 任務中的自我修正能力。
    \end{ZhAbstractDescription}
    
\end{ZhAbstract}
